\chapter{Security Hardening}
\label{ch:security}

Security is a primary design goal for Veilbox. The system is hardened at
every layer: kernel configuration reduces the attack surface, BusyBox
eliminates setuid binaries, Dropbear provides encrypted access, and
container runtimes use Linux security primitives.

\section{Kernel Security Features}

\subsection{Minimal Attack Surface}

The kernel is configured without unnecessary subsystems:
\begin{itemize}[nosep]
  \item No sound subsystem (CONFIG\_SOUND=n).
  \item No USB stack (CONFIG\_USB=n).
  \item No Wi-Fi or Bluetooth (CONFIG\_WLAN=n, CONFIG\_BT=n).
  \item No ATA or SCSI drivers (CONFIG\_ATA=n, CONFIG\_SCSI=n).
  \item No filesystems beyond ext4, proc, sysfs, tmpfs, devtmpfs,
        overlay, squashfs.
  \item No PCMCIA, FireWire, InfiniBand, or ISDN.
\end{itemize}

\subsection{Mandatory Security Options}

\begin{lstlisting}
CONFIG_SECCOMP=y           # Syscall filtering
CONFIG_SECCOMP_FILTER=y    # BPF-based filters
CONFIG_SECURITY=y          # Security framework
CONFIG_SECURITY_YAMA=y     # Ptrace scope restrictions
CONFIG_DEFAULT_SECURITY_YAMA=y
CONFIG_SECURITY_NETWORK=y  # Network security hooks
\end{lstlisting}

\subsection{CGroup Device Control}

CGroup device controllers restrict container access to host devices:

\begin{lstlisting}
CONFIG_CGROUP_DEVICE=y
\end{lstlisting}

\section{Userspace Security}

\subsection{No Setuid Binaries}

The Veilbox root filesystem contains no setuid binaries. BusyBox can be
configured with setuid support disabled:

\begin{lstlisting}[language=bash]
# BusyBox config: disable setuid
# CONFIG_FEATURE_SUID is not set
\end{lstlisting}

\subsection{No Package Manager}

Without a package manager, there is no mechanism for an attacker to
install software, exploit repository vulnerabilities, or run
post-install scripts.

\subsection{No Compilers}

No compilers or interpreters (other than busybox ash) are present in the
root filesystem. An attacker who gains shell access cannot compile
exploit code.

\subsection{Read-Only Root}

The root filesystem is embedded in the kernel binary. Runtime modification
of system binaries is impossible — changes are lost on reboot.

\subsection{Minimal Services}

Only two network services run: Dropbear (SSH) and containerd (local Unix
socket only). No HTTP, FTP, SMTP, NFS, SMB, or any other services.

\section{SSH Hardening}

\begin{itemize}[nosep]
  \item \textbf{ED25519 keys}: Modern, fast, secure host keys.
  \item \textbf{PBKDF2 password hash}: /etc/shadow uses SHA-512 rounds.
  \item \textbf{Port forwarding}: QEMU limits SSH to port forwarding
        only (no public guest IP by default).
  \item \textbf{Key-based auth}: Optional but supported with auto-generated
        ED25519 key pair.
\end{itemize}

\section{Container Security Primitives}

\subsection{Namespace Isolation}

Each container runs in isolated namespaces:
\begin{itemize}[nosep]
  \item \textbf{PID}: Container processes cannot see host processes.
  \item \textbf{Network}: Container has its own network stack.
  \item \textbf{Mount}: Container has its own filesystem hierarchy.
  \item \textbf{UTS}: Container has its own hostname.
  \item \textbf{IPC}: Container has its own IPC resources.
  \item \textbf{User}: (Optional) Container root can be mapped to an
        unprivileged host user.
\end{itemize}

\subsection{Seccomp}

Seccomp-BPF filters restrict the system calls available to container
processes. The default Docker/containerd seccomp profile allows
approximately 300 of the 400+ Linux syscalls, blocking dangerous ones
like \texttt{module\_init}, \texttt{kexec\_load}, and \texttt{ptrace}.

\subsection{Capability Dropping}

Linux capabilities are dropped from container processes:
\begin{lstlisting}
CAP_SYS_ADMIN     # Mount, namespace operations
CAP_NET_ADMIN     # Network configuration
CAP_SYS_MODULE    # Kernel module loading
CAP_SYS_RAWIO     # I/O port access
CAP_SYS_PTRACE    # Process tracing
\end{lstlisting}


\section{Container Security Architecture}

\subsection{Defense in Depth}

Veilbox implements multiple layers of container security:

\begin{enumerate}[nosep]
  \item \textbf{Namespaces}: Each container has isolated views of the
        system (processes, network, mounts, etc.).
  \item \textbf{Cgroups}: Resource limits prevent denial-of-service
        from a single container.
  \item \textbf{Seccomp-BPF}: System call filtering blocks dangerous
        syscalls.
  \item \textbf{Capability dropping}: Removes privileges like
        CAP\_SYS\_ADMIN and CAP\_NET\_ADMIN.
  \item \textbf{Read-only rootfs}: Containers can be started with
        read-only root filesystems.
  \item \textbf{No new privileges}: The NO\_NEW\_PRIVS flag prevents
        privilege escalation via setuid binaries.
\end{enumerate}

\subsection{Seccomp Default Profile}

The default seccomp profile blocks approximately 60 system calls that
are dangerous or unnecessary for containers:

\begin{lstlisting}
# Blocked syscalls (partial list)
kexec_load        # Loading a different kernel
open_by_handle_at # Bypass mount namespace
init_module       # Loading kernel modules
finit_module      # Loading kernel modules
delete_module     # Unloading kernel modules
ptrace            # Process tracing and debugging
swapon            # Enabling swap
swapoff           # Disabling swap
sethostname       # Changing hostname
setdomainname     # Changing domain name
\end{lstlisting}

\subsection{SELinux Considerations}

Veilbox does not use SELinux (it adds complexity and binary size). For
systems requiring SELinux:

\begin{lstlisting}
CONFIG_SECURITY_SELINUX=y
CONFIG_DEFAULT_SECURITY_SELINUX=y
\end{lstlisting}

However, this requires SELinux policies and tools which add approximately
10 MB to the initramfs.

\section{Audit Trail}

The system log captures security-relevant events:

\begin{lstlisting}
# Monitor authentication attempts
grep "dropbear" /var/log/messages

# Monitor container starts/stops
grep "containerd" /var/log/messages

# Monitor failed login attempts
grep "Failed" /var/log/messages
\end{lstlisting}

\section{Network Security}

\subsection{Firewall Rules}

Basic iptables rules for container isolation:

\begin{lstlisting}
# Default: drop incoming, allow outgoing
iptables -P INPUT DROP
iptables -P FORWARD DROP
iptables -P OUTPUT ACCEPT

# Allow loopback
iptables -A INPUT -i lo -j ACCEPT

# Allow established connections
iptables -A INPUT -m state --state ESTABLISHED,RELATED -j ACCEPT

# Allow SSH
iptables -A INPUT -p tcp --dport 22 -j ACCEPT

# Allow container traffic to external
iptables -t nat -A POSTROUTING -o eth0 -j MASQUERADE
\end{lstlisting}
